\section{Boundary geometry and the modulus criterion}
\label{sec:geometry}

We collect the geometric estimates used throughout the proof, then state
the classical uniformization input. Let $\Gamma$ be the Cayley graph of
$G$ for a finite symmetric generating set, with unit edge lengths. Write
$e$ for its identity vertex and $|v|=d_\Gamma(e,v)$.
The Gromov product is
\[
 (y|z)_o=\tfrac12\bigl(d_\Gamma(o,y)+d_\Gamma(o,z)-d_\Gamma(y,z)\bigr).
\]
Boundary products are defined by limits along rays, up to a uniformly
bounded additive error. A visual metric $d$ on $Z=\partial G$ satisfies
\begin{equation}
 C_v^{-1}a^{-(y|z)_e}\le d(y,z)\le C_va^{-(y|z)_e},
 \qquad y\ne z,
 \label{eq:visual}
\end{equation}
for some $a>1$. Set $r_k=a^{-k}$ for integers $k\ge0$.
We use the standard product inequality
$(y|z)_e\ge\min\{(y|w)_e,(w|z)_e\}-O(1)$ and its geometric form:
geodesics with product at least $j-C$ stay a bounded synchronous
distance apart through depth $j$, provided both have that length.
The bound depends on $C$ and the
hyperbolicity constant. These facts, visual metrics, and thin ideal
triangles are recalled in \cite[III.H.3]{BH}.

\subsection{Expansions}

Left multiplication by the inverse of a vertex far along a ray magnifies
a small boundary ball. The following formula also controls the complement
of that ball; this second feature will be essential.

\begin{lemma}[Expansion formula]
\label{lem:expansion}
Let $v$ be the depth-$k$ vertex on a based ray to $x\in Z$, put
$h=v^{-1}$ and $R=r_k$, and set $D_y=\max\{R,d(x,y)\}$.
There are constants $c_E,C_E>0$, independent of these choices, such that
\begin{equation}
 c_E\frac{R\,d(y,z)}{D_yD_z}
 \le d(hy,hz)\le C_E\frac{R\,d(y,z)}{D_yD_z}.
 \label{eq:expansion}
\end{equation}
In particular, on $B(x,CR)$ the map $h$ is bi-Lipschitz with scale
factor $R^{-1}$ and distortion depending only on $C$.
If $d(x,z_0)\ge\varepsilon_*>0$, $R<\varepsilon_*/2$, and $p=hz_0$, then
\begin{equation}
 d(hy,p)\lesssim R/D_y\quad(y\in Z),\qquad
 d(hy,p)\asymp R/D_y\quad(d(x,y)\le\varepsilon_*/2),
 \label{eq:puncture}
\end{equation}
and
\begin{equation}
 (v^{-1}|p)_e\ge k-C(\varepsilon_*).
 \label{eq:inverse-prefix}
\end{equation}
The constants in the last two estimates may depend on $\varepsilon_*$.
\end{lemma}
\begin{proof}
Base change and thin triangles give
\[
 (y|z)_v=(y|z)_e+k-(v|y)_e-(v|z)_e+O(1),
 \qquad (v|y)_e=\min\{k,(x|y)_e\}+O(1).
\]
Thus $a^{-(v|y)_e}\asymp D_y$. Exponentiating the first identity
and using $(hy|hz)_e=(y|z)_v+O(1)$ proves
\eqref{eq:expansion}. With $z=z_0$, the factor $D_{z_0}$ is bounded
above and below, and $d(y,z_0)$ is bounded above. When
$d(x,y)\le\varepsilon_*/2$, the latter is also bounded below.
This proves \eqref{eq:puncture}. Finally,
$(v^{-1}|v^{-1}z_0)_e=k-(v|z_0)_e+O(1)$, and
$(v|z_0)_e$ is bounded in terms of $\varepsilon_*$.
\end{proof}

\begin{lemma}[Prefixes]
\label{lem:prefix}
Every vertex $v$ lies a uniformly bounded distance from a depth-$|v|$
vertex on a based ray. With that ray's endpoint as $x$,
\eqref{eq:expansion} remains valid for $h=v^{-1}$ and $R=r_{|v|}$,
with uniform constants. Moreover:
\begin{enumerate}[label=\textup{(\roman*)}]
\item If $d(\xi,\eta)\le Cr_j$, depth-$j$ vertices on any based
rays to $\xi,\eta$ are a bounded distance apart, depending only on $C$.
\item If $|v|=k$ and $(v|\xi)_e\ge k-C$, every geodesic from $e$
to $v$ stays a bounded synchronous distance from a ray to $\xi$
through depth $k$.
\end{enumerate}
\end{lemma}
\begin{proof}
Translate a bi-infinite geodesic so that it passes through $v$.
The ideal triangle with third vertex $e$ puts $v$ uniformly near one
of its based sides; the nearby parameter differs from $|v|$ by a
bounded amount. A bounded change of basepoint changes products by a
bounded amount, preserving \eqref{eq:expansion}. The two remaining
assertions follow from \eqref{eq:visual} and the product estimate.
\end{proof}

Each fixed $g\in G$ acts bi-Lipschitzly on $Z$, with a constant
depending on $g$. The whole action is uniformly quasi-M\"obius:
base-change factors cancel in distance cross-ratios, giving a single
$C\ge1$ such that
\[
 C^{-1}\frac{d(x,z)d(y,w)}{d(x,w)d(y,z)}
 \le\frac{d(gx,gz)d(gy,gw)}{d(gx,gw)d(gy,gz)}
 \le C\frac{d(x,z)d(y,w)}{d(x,w)d(y,z)}
\]
for every distinct quadruple and every $g$.

\subsection{Packing and paths}

\begin{lemma}[Packing]
\label{lem:packing}
There are $Q\ge2$ and $C_P<\infty$ such that an $s$-separated subset
of a ball of radius $t$ has at most $C_P(t/s)^Q$ points whenever
$0<s\le t\le\diam Z$. In particular, $Z$ is doubling.
\end{lemma}
\begin{proof}
Choose $k$ with $r_k\asymp t$. Rays ending in the given ball have
their depth-$k$ vertices in a uniformly bounded graph ball, by
\Cref{lem:prefix}. At depth $k+j$ there are at most $CA^j$ possible
continuations, where $A>1$ depends on the valence. Endpoints whose
rays share such a vertex are at visual distance at most $C'r_{k+j}$.
Choose the least $j\ge0$ with $C'r_{k+j}<s$. The desired bound
follows with $Q\ge\max\{2,\log_a A\}$. Coarse scales are absorbed
by increasing the constant.
\end{proof}

\Needspace{5\baselineskip}
\begin{lemma}[Controlled paths]
\label{lem:paths}
There is $L\ge1$ such that any two points of $\overline B(x,s)$
join by a path in $B(x,Ls)$. Any two points of a continuum can also
be joined by a path in any prescribed open neighborhood of it.
\end{lemma}
\begin{proof}
For small $s$, expand at $x$ with $R\asymp s$ and choose $z_0$ a
uniformly positive distance from $x$; two fixed distinct points
suffice for this choice. By \eqref{eq:puncture}, the two image
points stay a uniform positive distance from $p=hz_0$.
Fix a homeomorphism $\phi:Z\to\Sphere$. Uniform continuity of
$\phi^{-1}$ places their spherical images a uniform distance from
$\phi(p)$. Join them outside a smaller round ball about $\phi(p)$.
Uniform continuity of $\phi$ makes the pullback path stay a uniform
$d$-distance from $p$. The upper bound in \eqref{eq:puncture}
puts its preimage under $h$ in $B(x,Ls)$. Increase $L$ for the
remaining scales, using path connectedness of $Z$.

For a continuum $C$ in an open set $V$, choose $\varepsilon>0$
with its $L\varepsilon$-neighborhood in $V$. Connectedness supplies
a finite $\varepsilon$-chain in $C$ between any two specified
points: the set reachable by such chains is relatively clopen.
Join consecutive points by the first assertion.
\end{proof}

\subsection{The classical reduction}

Choose a maximal $r_n$-separated set $T_n\subset Z$ and its cells
$B(q,2r_n)$, $q\in T_n$. These cover $Z$; fixed enlargements have
bounded overlap. For a path family $\mathcal F$, define
\[
 \Mod_2(\mathcal F,T_n)=\min_{\omega\ge0}
 \left\{\sum_{q\in T_n}\omega(q)^2:
 \sum_{B(q,2r_n)\cap\pi\ne\varnothing}\omega(q)\ge1
 \text{ for every }\pi\in\mathcal F\right\}.
\]
Only finitely many incidence vectors occur, so the minimum exists.
Fix a sufficiently small $d_0>0$ and put
\begin{equation}
 \mathcal F_0=\{\pi:\diam\pi\ge d_0\},\qquad
 M(n)=\Mod_2(\mathcal F_0,T_n).
 \label{eq:macro-modulus}
\end{equation}
We require $d_0<\diam Z/10$ and choose it so that every path of
diameter at least $R/2$ in $\overline B(x,3R)$ has image diameter
at least $d_0$ under a scale-$R$ expansion. This is possible by
\eqref{eq:expansion}. We also impose the smallness requirement in
\Cref{thm:uniformization} before fixing $d_0$ and $M(n)$.
The family is nonempty, and $0<M(n)<\infty$: all-one weights are
admissible, while any one path gives the lower bound $1/\#T_n$.

\begin{theorem}[Bonk--Kleiner; Bourdon--Kleiner]
\label{thm:uniformization}
If $\sup_n M(n)<\infty$, then $(Z,d)$ is quasi-M\"obius
homeomorphic to the round sphere.
\end{theorem}
\begin{proof}[Verification of the cited criterion]
Corollary 3.5 of \cite{BourdonKleiner}, derived from
\cite[Theorem 10.4]{BK}, applies to an arcwise connected,
approximately self-similar metric sphere. Approximate self-similarity
means that each open ball, with distances divided by its radius,
is uniformly bi-Lipschitz to an open subset of the space.
\Cref{lem:expansion} supplies this property; intermediate and large
radii change only the constant.

The criterion assumes bounded $2$-moduli of paths of diameter at
least a sufficiently small $d_0$ in successive graph approximations.
For a dyadic scale $2^{-m}$ choose $n$ with
$2^{-m}/a<r_n\le2^{-m}$. Our cells form a fixed $\kappa$-approximation
for any $\kappa>2a$: the inner balls of radius $2^{-m}/\kappa$
are disjoint, each is contained in its cell, and each cell is
contained in the radius-$\kappa2^{-m}$ ball about its center.
The modulus is exactly $M(n)$. Finitely many coarse scales have
finite modulus. No Coxeter-group assumption occurs in this criterion.
\end{proof}

\begin{proposition}[From the boundary to hyperbolic space]
\label{prop:action}
If $(Z,d)$ is quasi-M\"obius homeomorphic to the round sphere,
then $G$ admits a proper cocompact isometric action on $\HH^3$
with finite kernel.
\end{proposition}
\begin{proof}
Let $\psi:Z\to\Sphere$ be the uniformizing homeomorphism and write
$\mathcal G=\{\psi g\psi^{-1}:g\in G\}$. The boundary cross-ratio estimate
and the two fixed maps $\psi,\psi^{-1}$ give one quasi-M\"obius distortion
function for every element of $\mathcal G$. Fix $f\in\mathcal G$ and $x\in\Sphere$,
keep a fourth point away from $x$, and let two points at equal spherical
distance from $x$ approach $x$. In the cross-ratio inequality the factors
involving the fourth point tend to one, both before and after applying
$f$. Thus the infinitesimal linear distortion has a bound independent
of $f$ and $x$; the radius at which the estimate is accurate need not
be uniform in $f$. The metric characterization of planar
quasiconformality, applied in conformal charts, makes each $f$
quasiconformal. At its almost-everywhere nonsingular differentiability
points, the same infinitesimal bound controls the ratio of singular values
of $Df$, and hence its analytic dilatation. This gives a common
quasiconformal bound. For a reversing map apply the characterization
after a conformal reflection; see \cite[Theorems 2.5.3--2.5.4]{AIM}
and the metric characterization in \cite[Theorem 1.4 and Remark 1.8(b)]{HK}.

We recall explicitly the invariant-structure argument, including
orientation reversal. This is the invariant-conformal-structure method of
Sullivan, with the circumcenter construction detailed by Tukia
\cite[pp.~73--76]{Tukia}. Analytic quasiconformal maps have
almost-everywhere defined nonsingular weak differentials, and they and
their inverses preserve null sets \cite[Theorem 3.1.2]{AIM}.
Their differentials satisfy the almost-everywhere composition rule
\cite[Lemma 10 and Theorem 5]{AhlforsBers}.
Since $\mathcal G$ is countable, remove the $\mathcal G$-orbit
of the union of all exceptional sets for these assertions and for all
pairs of maps. Its complement $X$ is invariant and conull. At $x\in X$,
let $\mathcal C_x$ be the space of positive definite quadratic forms on
$T_x\Sphere$ modulo positive scalars. This is a hyperbolic plane, and
pullback by any invertible linear map, of either orientation, is an
isometry between the corresponding fibers
\cite[Section 2.3, equations (2.17)--(2.18) and Theorem 2.3.2]{AIM}. The orbit
$\{f^*q_{\rm round}(x):f\in\mathcal G\}\subset\mathcal C_x$ is uniformly
bounded, by the common quasiconformal bound, and contains the round
form. Its unique hyperbolic circumcenter defines a uniformly bounded
measurable conformal structure $q(x)$. Indeed, circumcenters of the
increasing finite subsets from an enumeration are measurable and
converge to that circumcenter: their radii increase, their centers stay
in a fixed compact ball, and every subsequential limit is the unique
center of the whole orbit. The composition rule permutes this orbit
under pullback, so $f^*q=q$ on $X$ for every $f\in\mathcal G$. Extend $q$
arbitrarily across the null complement of $X$.

The measurable Riemann mapping theorem
\cite[Theorem 6]{AhlforsBers} gives a quasiconformal
homeomorphism $F:\Sphere\to\Sphere$ with $F^*q_{\rm round}=q$ almost
everywhere. Invariance and the analytic composition rule show that
every $FfF^{-1}$ is conformal or anticonformal, according to the
orientation of $f$. Such sphere homeomorphisms are M\"obius or
anti-M\"obius. Their unique extensions to hyperbolic isometries therefore
define a homomorphism $\rho:G\to\Isom(\HH^3)$.

For completeness, properness and cocompactness follow directly
from the action on ordered distinct triples $\mathcal T(Z)$.
On any compact set of triples all pairwise distances are bounded
below. For $|v|\to\infty$, at least two points of such a triple
stay away from the expansion center for $v^{-1}$; their image
distance is $O(r_{|v|})$. Thus $v^{-1}K\cap K=\varnothing$
eventually for every compact $K\subset\mathcal T(Z)$.
Conversely, expand at one point of a closest pair with scale
comparable to their distance. Formula \eqref{eq:expansion} makes
all three image distances uniformly positive: if the third point
is nearby, all denominator factors have that scale; otherwise
its distances to the pair cancel its denominator factor.
Triples already uniformly separated need no expansion.
Hence the action on triples is proper and cocompact.

These properties survive boundary conjugacy. Let $I=\Isom(\HH^3)$.
The stabilizer $H$ of a standard ordered boundary triple has two
elements: the identity and reflection in its hyperbolic plane.
Choose $o\in\HH^3$ fixed by $H$, and let $K_o$ be its compact
stabilizer in $I$. The equivariant map
\[
 \mathcal T(\Sphere)\cong I/H\longrightarrow I/K_o\cong\HH^3
\]
is continuous, onto, and proper: isometries sending $o$ into a
compact set form a compact set. Properness and cocompactness
therefore descend. The kernel fixes $o$, so is finite.
\end{proof}
